\begin{description}
\item[(A)] First evaluate the radius of the geostationary orbit using the sidereal period
$T=23^{\mathrm h}56^{\mathrm m}4^{\mathrm s}$,
\[
T^2=\frac{4\pi^2R_S^3}{GM_E}\quad\Longrightarrow\quad R_S=4.2163\times10^7\,\mathrm{m}=42163\,\mathrm{km}.
\]
\ioaafig{2023_TH_Q12-s1.pdf}
In the figure, E is the center of Earth, S the satellite, O the observer on Earth, and Z the observer's zenith.
The latitude is $\varphi$ and the elevation angle is $\upsilon$.
$EO=R_E$ is Earth's radius, $ES=R_S$ the satellite's orbital radius, and $OS=x$ the observer--satellite distance.
The limits of the visible latitude range occur when $\upsilon=0$. The sine law gives
\[
\frac{\sin\widehat{EOS}}{R_S}=\frac{\sin\widehat{OSE}}{R_E},
\qquad
\frac{\sin90^\circ}{R_S}=\frac{\sin(180^\circ-(90^\circ+\varphi))}{R_E}.
\]
Hence $\varphi=\cos^{-1}(R_E/R_S)\simeq81.3^\circ$.
The maximum latitudes are $81.3^\circ$ N and $81.3^\circ$ S. \hfill\mbox{(12 points)}
\item[(B)] Find the location where $\widehat{OSE}=5^\circ$ and $\upsilon=0$.
The triangle SOE is right-angled, with $\widehat{EOS}=90^\circ$. Therefore,
\[
\sin\widehat{OSE}=\frac{R_E}{R_S}=\sin5^\circ=0.0872,
\qquad R_S=\frac{6378}{0.0872}=73179\ \mathrm{km}.
\]
\hfill\mbox{(10 points)}
\item[(C)] The total energy is
\[
E_{\mathrm{tot}}=\tfrac12mv^2-G\frac{M_Em}{r}.
\]
The centripetal acceleration relation $v^2/r=GM/r^2$ gives
\[
E_{\mathrm{tot}}=-\frac12G\frac{M_Em}{r}.
\]
Substitution at the two locations gives $E_{\mathrm{tot},1}=-4.727\times10^9$ J and
$E_{\mathrm{tot},2}=-2.723\times10^9$ J. The energy has increased by $2.00\times10^9$ J. \hfill\mbox{(10 points)}
\item[(D)] Gravitational acceleration is inversely proportional to the square of the radius,
$g=GM/r^2$. Thus
\[
\frac{g_E}{g_B}=\left(\frac{r_B}{r_E}\right)^2
=\left(\frac{73179000}{6378000}\right)^2=131.6.
\]
\hfill\mbox{(8 points)}
\end{description}
